%% ma: L10-B14-0069
%% lop: 10
%% chuong: 5
%% bai: 14
%% dang: 10.14.7
%% loai: TN4
%% muc: TH
%% dap_an: "B"
%% nguon: "chuyen-de-cac-so-dac-trung-cua-mau-so-lieu-khong-ghep-nhom-toan-10-chuong-trinh-moi.pdf (GV. Phan Nhật Linh) (P05), Bài 03 Dạng 2 Phần I Câu 13"
%% kien_thuc: "Lớp 10 Bài 13 (số trung bình), Bài 14 (phương sai, độ lệch chuẩn)"
%% trang_thai: chinh-thuc
%% ghi_chu: "Tự duyệt theo yêu cầu giáo viên 22/10/2026."
\begin{ex}
Cho mẫu số liệu $1$; $2$; $3$; $4$; $5$; $6$; $7$; $8$. Độ lệch chuẩn của mẫu số liệu gần bằng
\choice{$3{,}30$}{\True $2{,}30$}{$4{,}30$}{$5{,}30$}
\loigiai{$\overline x=4{,}5$; $s^2=\dfrac{2(3{,}5^2+2{,}5^2+1{,}5^2+0{,}5^2)}8=\dfrac{21}4$; $s=\sqrt{\dfrac{21}4}\approx2{,}29$, gần $2{,}30$ nhất. Chọn B.}
\end{ex}
